% ---------------------------------------------------------------------
% EG author guidelines plus sample file for EG publication using LaTeX2e input
% D.Fellner, v1.17, Sep 23, 2010


\title[Peridynamics-Based Fracture Animation for Elasto-Plastic Solids]%
      {Peridynamics-Based Fracture Animation for Elasto-Plastic Solids}

% for anonymous conference submission please enter your SUBMISSION ID
% instead of the author's name (and leave the affiliation blank) !!
\author[paperID]{paperID}

%\author[D. Fellner \& S. Behnke]
%       {D.\,W. Fellner\thanks{Chairman Eurographics Publications Board}$^{1,2}$
%        and S. Behnke$^{2}$
%%        S. Spencer$^2$\thanks{Chairman Siggraph Publications Board}
%        \\
%% For Computer Graphics Forum: Please use the abbreviation of your first name.
%         $^1$TU Darmstadt \& Fraunhofer IGD, Germany\\
%         $^2$Institut f{\"u}r ComputerGraphik \& Wissensvisualisierung, TU Graz, Austria
%%        $^2$ Another Department to illustrate the use in papers from authors
%%             with different affiliations
%       }

% ------------------------------------------------------------------------

% if the Editors-in-Chief have given you the data, you may uncomment
% the following five lines and insert it here
%
% \volume{27}   % the volume in which the issue will be published;
% \issue{1}     % the issue number of the publication
% \pStartPage{1}      % set starting page


%-------------------------------------------------------------------------
\begin{document}

% \teaser{
%  \includegraphics[width=\linewidth]{eg_new}
%  \centering
%   \caption{New EG Logo}
% \label{fig:teaser}
% }

\maketitle

\begin{abstract}
Key sale point:
the first peridynamics framework in graphics to simulate deformation and fracture of elasto-plastic materials.
\begin{classification} % according to http://www.acm.org/class/1998/
\CCScat{Computer Graphics}{I.3.3}{Picture/Image Generation}{Line and curve generation}
\end{classification}

\end{abstract}





%-------------------------------------------------------------------------
\section{Introduction}

%-------------------------------------------------------------------------
\section{Related Work}

%-------------------------------------------------------------------------
\section{Background}

%-------------------------------------------------------------------------
\section{Overview}

%-------------------------------------------------------------------------
\section{Constitutive Model}

%-------------------------------------------------------------------------
\section{Fracture}

%-------------------------------------------------------------------------
\section{Implementation}

%-------------------------------------------------------------------------
\section{Results}

%-------------------------------------------------------------------------
\section{Discussion}

%-------------------------------------------------------------------------
%\bibliographystyle{eg-alpha}
%\bibliographystyle{eg-alpha-doi}

%\bibliography{references}

%-------------------------------------------------------------------------
%\newpage

\end{document}
